\documentclass[12pt,a4paper,twoside]{thesis}

%% PACKAGES
\usepackage[utf8]{inputenc}
\usepackage[T1]{fontenc}
\usepackage{algorithm}
\usepackage[noend]{algpseudocode}
\usepackage{amsmath}
\usepackage{amsthm}
\usepackage{amsfonts}
\usepackage{amssymb}
\usepackage{amsthm}
\usepackage{bytefield}
\usepackage{cancel}
\usepackage[dvips]{color}
\usepackage{enumerate}
\usepackage{epigraph}
\usepackage{fancyhdr}
\usepackage{graphicx}
\usepackage{hyperref}
\usepackage{indentfirst}
\usepackage{mathtools}
\usepackage{minted}
\usepackage{makeidx,shortvrb,latexsym}
\usepackage{supertabular}
\usepackage{tikz}

%% setting epigraphs
\renewcommand{\epigraphsize}{\small}
\setlength{\epigraphwidth}{0.8\textwidth}
\let\origepigraph\epigraph
\renewcommand\epigraph[2]{\origepigraph{\textit{#1}}{\textsc{#2}}}

%% bytefield shit
\newcommand{\colorbitbox}[3]{%
\rlap{\bitbox{#2}{\color{#1}\rule{\width}{\height}}}%
\bitbox{#2}{#3}}


%% COMMANDS
\DeclarePairedDelimiter{\floor}{\lfloor}{\rfloor}
\DeclarePairedDelimiter{\ceil}{\lceil}{\rceil}
\DeclarePairedDelimiter{\angular}{\langle}{\rangle}

\newcommand{\naturalN}{\mathbb{N}}
\newcommand{\naturalPrime}{\mathbb{P}}
\newcommand{\integerZ}{\mathbb{Z}}
\newcommand{\bigO}[1]{\ensuremath{\operatorname{O}\left(#1\right)}}
\newcommand{\openssl}{\textsc{OpenSSL}\ }
%% here adopting Wikipedia's notation <https://en.wikipedia.org/wiki/Isqrt>
\newcommand{\dsqrt}[1]{\ensuremath{isqrt(#1)}}
\newcommand{\idiv}{\ensuremath{//}}
\newcommand{\strong}[1]{\textbf{#1}}
\newcommand{\eulerphi}[1]{\varphi(#1)}
\newcommand{\abs}[1]{\left|#1\right|}
\newcommand{\rfrac}[2]{{}^{#1}\!/_{#2}}
\newcommand{\getsRandom}{\xleftarrow{r}}


\theoremstyle{plain}
\newtheorem*{theorem*}{Theorem}
\newtheorem*{definition*}{Definition}

\makeindex
\let\origdoublepage\cleardoublepage
\newcommand{\clearemptydoublepage}{%
\clearpage
{\pagestyle{empty}\origdoublepage}%
}
\let\cleardoublepage\clearemptydoublepage
% Note that colour boxes around links are NOT printed.
% The text itself can be coloured, replacing the bounding box, but
% when printing this may appear illegible;
% link colors can be set to black for printing purposes, like so:
%\hypersetup{colorlinks,%
%	citecolor=black,%
%	filecolor=black,%
%	linkcolor=black,%
%	urlcolor=black}
% Also note that this will conflict with the color package called
% earlier in this document if that is not set to the right option (e.g.
% dvips or pdflatex).
\begin{document}
%\phd  %se e' la tesi di dottorato, altrimenti non mettere nulla

\university{Universit\`a degli Studi di Trento}
\faculty{Facolt\`a di Scienze Matematiche Fisiche e Naturali}
\dept{Dipartimento di Scienze Informatiche}
\Logo{logo_unitn.jpg}
%% titolo del dottorato
% \phdtitle{}
%% titolo della tesi
\title{Question Authority}
\subtitle{An Inquiry into The Secure Layer}

\author{Michele Orr\`u}
\supervisor{Prof. Massimiliano Sala}
%% Uncomment the following two lines if a co-relator is present.
%\twosupervisors
%\firstreader{Prof. Ettore Miron}

% capo della scuola di dottorato o controrelatore
\secondreader{Prof. Nara Stabbocchi}
\accademico{Anno accademico $2013/2014$}

\frontespizio     % questo e` il frontespizio esterno, cioe' senza firme
\cleardoublepage
\signaturepage    % questo e` il frontesizio interno con le firme

%% DEDICA
%non e' ovviamente obbligatoria


\cleardoublepage \setcounter{page}{1} \pagenumbering{roman}
\pagestyle{plain} \tableofcontents
%\listoffigures
%\addcontentsline{toc}{chapter}{Elenco delle figure}
%\renewcommand{\listalgorithmname}{Elenco degli algoritmi}
%\listofalgorithms
%\addcontentsline{toc}{chapter}{Elenco degli algoritmi}
%\addcontentsline{toc}{chapter}{Introduction}
%% INTRODUCTION
%\include{ack}
%\addcontentsline{toc}{chapter}{Acknowledgment}
%\cleardoublepage
%\include{Introduction}
%\addcontentsline{toc}{chapter}{Introduction}
\cleardoublepage

%% THESIS BODY
\pagestyle{fancy} \pagenumbering{arabic} \mainmatter
\vspace*{3in}
\epigraph{
    Many persons who are not conversant with mathematical studies imagine that
    because the business of [Babbage's Analytical Engine] is to give results in
    numerical notation, the nature of the processes must consequently be
    arithmetical and numerical, rather than algebraical and analytical. This is an
    error. The engine can arrange and combine its numerical quantities exactly as if
    they were letters or any other general symbols; and in fact it might bring out
    its results in algebraical notation, were provisions made accordingly.}
         {Augusta Ada, Countess of Lovelace}


%% there is no copyright, but the right to copy.
\begin{figure}[b]
  \centering
  \includegraphics[width=80pt]{kopimi.png}
\end{figure}


\chapter{Preface}

Even if the basic RSA key generation algorithm is fairly straightforward, it
turns out that any software willing to provide such a feature does have to test
the candidate key against a substantious number of threats before claiming its
security.

The purpose of this project is to examine the TLS protocol, to study in depth
the \openssl library, and to survey some of the attacks to which a bad key
generation is exposed. On the footprint of ~\cite{20years}, where most of these
attacks have already been analyzed, we are going to describe the mathematical
basis of each attack, reason about a possible solution in procedural
programming, and finally, give some hints about a distributed version of it.

Besides the pseudocode already available in this document, the project led to the
development of a real, open, C implementation consultable at
\small{\url{https://github.com/mmaker/bachelor}}.

%%% Local Variables:
%%% mode: latex
%%% TeX-master: "question_authority"
%%% End:


\part{Prequisites}
\chapter{The Secure Socket Layer \label{chap:ssl}}


cos'e
differenze tra le varie versioni
la certification autority
%% in homage to sir Isacc Newton
\chapter{Mathematical Principles\label{chap:preq}}

In this chapter we formalize the notation used in the rest of the thesis, and
furthermore attempt to discuss and study the elementary functions on which the
project has been grounded.
\\
The $\ll$ and $\gg$ are respectively used with the meaning of left and right
bitwise shift, as usual in computer science.
\\
The $\dsqrt$ function will be defined in section \ref{sec:preq:sqrt}, with the
acceptation of discrete square root.
\\
The logarithmic $\log$ function is assumed to be in base two, i.e. $\log_2$.
\\
The $\idiv$ symbol is the integer division over $\naturalN$, i.e.
$a \idiv b = \floor{\frac{a}{b}}$, as usual in the python language.
\\
$\naturalPrime \subset \naturalN$ is the set containing all prime intgers.
\\
The binary operator $\getsRandom$, always written as $x \getsRandom S$, has the
meaning of ``pick a uniformly distributed random element $x$ from the set $S$''.
% XXX.  following Dan Boneh notation
\\
The summation in $\mathbb{F}_2$ is always expressed with the circled plus,
i.e. $a \xor b$.
%% Since it is equivalent to the bitwise xor, we are going to use
%% it as well in the pseudocode with the latter meaning.


%%\section{Number Theory}

%%What follows here is the definition and the formalization of some intuictive
%%concepts that later are going to be taken as granted:
%%the infinite cardinality of $\naturalPrime$,
%%the definition of \emph{smoothness}, and
%%the distribution of prime numbers in $\naturalN$.

\begin{definition*}[Smoothness]
  A number $n$ is said to be $\factorBase$-smooth if and only if all its prime
  factors are contained in $\factorBase$.
\end{definition*}

\begin{definition*}[Quadratic Residue]
  An integer $a$ is said to be a \emph{quadratic residue} $\mod n$ if it is
  congruent to a perfect square $\!\mod n$:
  \begin{equation*}
    x^2 = a \pmod{n}
  \end{equation*}
\end{definition*}

\begin{definition*}[Legendre Symbol]
  The \emph{Legendre Symbol}, often contracted as $\legendre{a}{p}$ is a
  function of two integers $a$ and $p$ defined as follows:
  \begin{equation*}
    \legendre{a}{p} = \begin{cases}
      0  & \text{if $a \equiv 0 \pmod{p}$} \\
      1  & \text{if $a$ is a quadratic residue modulo $p$} \\
      -1 & \text{if $a$ is a non-residue modulo $p$} \\
      \end{cases}
  \end{equation*}
\end{definition*}
\vfill


\section{Algorithmic Complexity Notation}
The notation used to describe asymptotic complexity follows the $\mathcal{O}$-notation,
abused under the conventions and limits of MIT's Introduction to Algorithms
\cite{MITalg}.

Let \bigO{g} be the asymptotic upper bound of g:
$$
\bigO{g(n)} = \{ f(n) : \exists n_0, c \in \naturalN \mid 0 \leq f(n) \leq cg(n)
\ \forall n > n_0 \}
$$

With $f(n) = \bigO{g(n)}$ we actually mean
$f(n) \in \bigO{g(n)}$.
Moreover, since the the expression ``running time'' has achieved a certain
vogue, we shall sometimes use this term as interchangeable with ``complexity'',
even though imprecise (\cite{Crandall} \S 1.1.4).

\section{Euclid's Greatest Common Divisor \label{sec:preq:gcd}}

Being the greatest common divisor a foundamental algebraic operation in the TLS
protocol, \openssl implemented it with the following signature:

\begin{minted}[fontsize=\small]{c}
  int BN_gcd(BIGNUM *r, BIGNUM *a, BIGNUM *b, BN_CTX *ctx);
\end{minted}

The computation proceeds under the well-known Euclidean algorithm, specifically
the binary variant developed by Josef Stein in 1961 \cite{AOCPv2}. This variant
exploits some interesting properties of $gcd(a, b)$:
\begin{enumerate}[(a)]
  \setlength{\itemsep}{1pt}
  \item if $a,\ b$ are even, then $gcd(a, b) = 2gcd(a/2, b/2)$;
  \item if $a$ is even and $b$ is odd, then $gcd(a, b) = gcd(a/2, b)$;
  \item  $gcd(a, b) = gcd(a-b, b)$, as in the standard Euclid algorithm;
  \item the sum of two odd numbers is always even.
\end{enumerate}

% Donald Knuth, TAOCP, "a binary method", p. 388 VOL 2
Both \cite{AOCPv2} and \cite{MITalg} analyze the running time of the
algorithm; \cite{MITalg}'s proof is fairly simpler and proceeds %elegantly
by induction.
Anyway, both show that algorithm ~\ref{alg:gcd} belongs to the class
\bigO{\log b}.

\begin{algorithm}[H]
  \caption{\openssl's GCD \label{alg:gcd}}
  \begin{algorithmic}[1]
    \Function{gcd}{$a, b$}
    \State $k \gets 0$
    \While{$b \neq 0$}
      \If{$a$ is odd}
        \If{$b$ is odd}
          \Comment by property (c) and (d)
          \State $a \gets (a-b) \gg 1$
        \Else
          \Comment by property (b)
          \State $b \gets b \gg 1$
        \EndIf
        \If{$a < b$} $a, b \gets b, a$ \EndIf

      \Else
        \If{$b$ is odd}
        \Comment by property (b)
          \State $a \gets a \gg 1$
          \If{$a < b$} $a, b \gets b, a$ \EndIf
        \Else
        \Comment by property (a)
          \State $k \gets k+1$
          \State $a, b \gets a \gg 1, b \gg 1$
        \EndIf
      \EndIf
    \EndWhile
    \State \Return $a \ll k$
    \EndFunction
  \end{algorithmic}
\end{algorithm}

Unfortunately, there is yet no known parallel solution that significantly improves
Euclid's \textsc{gcd}.


\section{Square Root \label{sec:preq:sqrt}}

Computing the square root is another important building block of the project,
though not available in \openssl\!.
Apparently,
% \openssl is a great pile of crap, as phk states
\openssl does only provide the discrete square root implementation using the
Tonelli/Shanks algorithm, which specifically solves in $x$ the equation
$x^2 = a \pmod{p}$, with $p \in \naturalPrime$:

\begin{minted}{c}
  BIGNUM* BN_mod_sqrt(BIGNUM* x, const BIGNUM* a, const BIGNUM* p,
                      const BN_CTX* ctx);
\end{minted}

Instead, we are interested in finding the the pair
$\angular{x, r} \in \naturalN^2 \mid x^2 + r = n$, that is, the integer part of
the square root of a natural number and its rest.
Hence, we did come out with our specific implementation, first using Bombelli's
algorithm, and later with the one of Dijkstra. Both are going to be discussed
below.

Unless otherwise specified, in the later pages we use $\sqrt{n}$ with the
usual meaning ``the half power of $n$'', while with $x, r = \dsqrt{n}$ we mean
the pair just defined.

\paragraph{Bombelli's Algorithm \label{par:preq:sqrt:bombelli}} dates back to
the XVI century, and approaches the problem of finding the square root by using
continued fractions. Unfortunately, we weren't able to fully assert the
correctness of the algorithm, since the original document
~\cite{bombelli:algebra} presents a difficult, inconvenient notation. Though,
for completeness' sake, we report in table
~\ref{alg:sqrt:bombelli} the pseudocode adopted and tested for its correctness.

\begin{algorithm}[H]
  \caption{Square Root: Bombelli's algorithm}
  \label{alg:sqrt:bombelli}
  \begin{algorithmic}[1]
    \Function{sqrt}{$n$}

    \State $i \gets 0; \quad g \gets \{\}$
    \While{$n > 0$}
    \Comment take pairs of digits and store them in $g$
      \State $g_i \gets n \pmod{100}$
      \State $n \gets n // 100$
      \State $i \gets i + 1$
    \EndWhile

    \State $x \gets 0; \quad r \gets 0$
    \For{$j =  i-1 \strong{ downto } 0$}
      \State $r = 100r + g_i$
      \Comment take next pair
      \For{$d = 0 \strong{ to } 9$}
      \Comment find gratest multiplier $d$
        \State $y' \gets d(20x + d)$
        \If{$y' > r$}  \textbf{break}
        \Else  \ \ $y \gets y'$
        \EndIf
      \EndFor
      \State $r \gets r - y$
      \State $x \gets 10x + d - 1$
      \Comment $d$ is the next digit
    \EndFor

    \State \Return $x, r$

    \EndFunction
  \end{algorithmic}
\end{algorithm}

For each digit of the result, we perform a subtraction, and a limited number of
multiplications. This means that the complexity of this solutions belongs to
\bigO{\log n \log n} = \bigO{\log^2 n}.

\begin{remark}
  Note that Bombelli actually has found a solution in $x$ for a slightly
  different equation than the one we initially formulated. Specifically, he
  found the pair $\angular{x, r}$ such that $(x+r)^2=a$, where $x$ is the mantissa,
  while $r$ is the decimal part. For our purpose this change is irrelevant: we
  just need to be able to distinguish perfect squares, and thus assert that $r$
  is zero.
\end{remark}

\paragraph{Dijkstra's Algorithm \label{par:preq:sqrt:dijkstra}} can be found in
\cite{Dijkstra:adop}, \S 8, p.61. There, Dijkstra presents an elightning
process for the computation of the square root, making only use of binary shift
and algebraic additions.
Specifically, the problem attempts to find, given a natual $n$, the integer $a$
that establishes:
\begin{align}
  \label{eq:preq:dijkstra_problem}
  a^2 \leq n \: \land \: (a+1)^2 > n
\end{align}

Take now the pair $\angular{a=0, b=n+1}$, and consider the inverval
$[a, b[$. We would like to reduce the distance between the upper bound $b$ and
the lower bound $a$, while holding the guard \ref{eq:preq:dijkstra_problem}:

\begin{align*}
  a^2 \leq n \: \land \: b > n
\end{align*}

%% XXX. I am not so sure about this, pure fantasy.
The speed of convergence is determined by the choice of the distance $d$, which
analougously to the dicotomic search problem, is optimal when
$d = (b-a) \idiv 2$.

\begin{algorithm}[H]
  \caption{Square Root: an intuitive, na\"ive implementation}
  \label{alg:sqrt:dijkstra_naif}
  \begin{algorithmic}[1]
    \Function{sqrt}{$n$}
    \State $a \gets 0; \quad b \gets n+1$
    \While{$a+1 \neq b$}
      \State $d \gets (b-a) \idiv 2$
      \If{$(a+d)^2 \leq n$} $a \gets a+d$
      \Comment increment left bound
      \ElsIf{$(b-d)^2 > n$} $b \gets b-d$
      \Comment increment right bound
      \EndIf
    \EndWhile
    \State \Return $a, a^2-n$
    \EndFunction
  \end{algorithmic}
\end{algorithm}
% heh, there's not much to explain here, that's almost the same in Dijkstra's
% book, excluding the inspirative familiar portrait that led to the insight of
% this change of varaibles.
Now optimization proceeds with the following change of variables:
\begin{enumerate}[a)]
  \setlength{\itemsep}{1pt}
  \setlength{\parskip}{0pt}
  \setlength{\parsep}{0pt}
\item $c = b-a$,
\item $p = ac$,
\item $q = c^2$,
\item $r = n-a^2$;
\end{enumerate}
resulting into algorithm \ref{alg:sqrt:dijkstra}.
For any further details, the reference is still \cite{Dijkstra:adop}.

\begin{algorithm}[H]
  \caption{Square Root: final version}
  \label{alg:sqrt:dijkstra}
  \begin{algorithmic}[1]
    \Function{sqrt}{$n$}
    \State $p \gets 0; \quad q \gets 1; \quad r \gets n$
    \While{$q \leq n$} $q \gets q \ll 2$ \EndWhile
    \While{$q \neq 1$}
      \State $q \gets q \gg 2$
      \State $h \gets p+q$
      \State $p \gets q \ll 1$
      \State $h \gets 2p + q$
      \If{$r \geq h$}
        \State $p \gets p+q$
        \State $r \gets r-h$ \EndIf
    \EndWhile
    \State \Return $p, r$
    \EndFunction
  \end{algorithmic}
\end{algorithm}

A fair approximation of the magnitude of the Dijkstra algorithm can be studied
by looking at the pseudocode in ~\ref{alg:sqrt:dijkstra_naif}. Exactly as in
the dicotomic search case, we split the interval $[a, b]$ in half on each step,
and choose whether to take the leftmost or the rightmost part. This results in
$log(n+1)$ steps. During each iteration, instead, as we have seen in
~\ref{alg:sqrt:dijkstra} we just apply summations and binary shifts, which are
upper bounded by \bigO{\log{n}/2}. Thus, the order of magnitude belongs to
\bigO{\log^2{n}}.

\paragraph{}
Even if both algorithms presented have \emph{asymptotically} the same
complexity, we believe that adopting the one of Dijkstra has lead to a
pragmatic, substantial performance improvement.


\section{The RSA Cipher \label{sec:preq:rsa}}

The RSA cryptosystem, invented by Ron Rivest, Adi Shamir, and Len Adleman
~\cite{rsa}, was first published in August 1977's issue of
\emph{Scientific American}. In its basic version, this \emph{asymmetric} cipher
works as follows:
\begin{itemize}
  \item choose a pair $\angular{p, q}$ of \emph{random} \emph{prime} numbers;
    let $N$ be the product of the two, $N=pq$, and call it \emph{public modulus};
  \item choose a pair $\angular{e, d}$ of \emph{random} numbers, both in
    $\integerZ^*_{\varphi(N)}$, such that one is the multiplicative inverse of the
    other, $ed \equiv 1 \pmod{\varphi(N)}$ and $\varphi(N)$ is Euler's totient
    function;
\end{itemize}
Now, call $\angular{N, e}$ \emph{public key}, and $\angular{N, d}$
\emph{private key}, and let the encryption function $E(m)$ be the $e$-th power of
the message $m$:
\begin{align}
  \label{eq:rsa:encrypt}
  E(m) = m^e \pmod{N}
\end{align}
while the decryption function $D(c)$ is the $d$-th power of the ciphertext $c$:
\begin{align}
  \label{eq:rsa:decrypt}
  D(c) = c^d \equiv E(m)^d \equiv m^{ed} \equiv m \pmod{N}
\end{align}
that, due to Fermat's little theorem, is the inverse of $E$.

\paragraph{}
%% less unless <https://www.youtube.com/watch?v=XnbnuY7Kxhc>
From now on, unless otherwise specified, the variable $N=pq$ will always refer
to the public modulus of a generic RSA keypair, with
$p, q$ being the two primes factorizing it, such that $p > q$.
 Again, $e, d$ will respectively refer to the public
exponent and the private exponent.

%%% Local Variables:
%%% mode: latex
%%% TeX-master: "question_authority"
%%% End:


\part{Factorization Methods}
\chapter{Fermat's factorization method \label{chap:fermat}}

Excluding the trial division, Fermat's method is the oldest known systematic
method for factorizing integers. Even if its algorithmic complexity is not
among the most efficient, it holds still a practical interest whenever
the two primes are sufficiently close.
Indeed, \cite{DSS2009} \S B.3.6 explicitly recommends that
$|p-q| \geq \sqrt{N}2^{-100}$
for any key of bitlength $1024,\ 2048,\ 3072$ in order to address this kind of
threat.\\
%% it would be nice here to explain that this magic 2^100 is just about wonting
%% the most significant digits to be different.
The basic idea is to attempt to write $N$ as a difference of squares,
\begin{align}
\label{eq:fermat_problem}
x^2 - N = y^2
\end{align}

So, we start by $x = \ceil{\sqrt{N}}$ and check that $x^2-N$ is a perfect
square. If it isn't, we iteratively increment $x$ and check again, until we
find a pair $\angular{x, y}$ satisfying equation \ref{eq:fermat_problem}.
Once found, we claim that $N = pq = (x+y)(x-y)$; it is indeed true that, if we
decompose $x^2 - y^2$ as difference of squares, then it is immediately clear
that $x+y \mid N \ \land \  x-y \mid N$, and that both are non-trivial
divisors.

\paragraph{Complexity} \cite{riesel} contains a detailed proof for the
complexity of this algorithm, which is
$\bigO{\frac{(1-k)^2}{2k} \sqrt{N}} \;\;,  0 < k < 1$. We summarize it down
below here to better clarify the limits of this algorithm.

\begin{proof}
  Since, once we reach the final step $x_f$ it holds $N = pq = x_f^2 - y_f^2$,
  the number of steps required to reach the result is:
  \begin{align*}
    x_f - \sqrt{N} &= \frac{p + q}{2} - \sqrt{N} \\
                   &= \frac{p + \frac{N}{p}}{2} - \sqrt{N} \\
                   &= \frac{(\sqrt{N} - p)^2}{2p}
  \end{align*}
  If we finally suppose that $p = k\sqrt{N}, \; 0 < k < 1$, then the number of cycles
  becomes
  $\frac{(1-k)^2}{2k} \sqrt{N}$.
\end{proof}

\begin{remark}
  Note that, for the algorithm to be effective, the two primes must be
  ``really close'' to $\sqrt{N}$. As much as the lowest prime gets near to
  $1$, the ratio $\frac{(1-k)^2}{2k}$ becomes larger, until the actual magnitude
  of this factorization method approaches \bigO{N}.
\end{remark}

\section{An Implementation Perspective \label{sec:fermat:implementation}}

At each iteration, the $i-$th state is hold by the pair $\angular{x, x^2}$.\\
The later step, described by $\angular{x+1, (x+1)^2}$ can be computed efficiently
considering the square of a binomial: $\angular{x+1, x^2 + 2x + 1}$.
The upper-bound, instead, is reached when
$ \Delta = p - q  = x + y - x + y = 2y > 2^{-100}\sqrt{N}$.

Algorithm ~\ref{alg:fermat} presents a simple implementation of this
factorization method, taking into account the small optimizations
aforementioned.

\begin{algorithm}[H]
  \caption{Fermat Factorization \label{alg:fermat}}
  \begin{algorithmic}[1]
    \Function{fermat}{\PKArg}
    \State $x \gets \floor{\sqrt{N}}$
    \State $x' \gets x \cdot x$

    \Repeat
    \State $x' \gets x' + 2x + 1$
    \State $x \gets x+1$
    \State $y, rest \gets \dsqrt{x' - N}$
    \Until{ $rest \neq 0 \strong{ and } y < \frac{\sqrt{N}}{2^{101}}$ }
    \Comment i.e., \ref{eq:fermat_problem} holds?

    \If{ $rest = 0$ }
    \State $p \gets x+y$
    \State $q \gets x-y$
    \State \Return $p, q$
    \Else
    \State \Return \textbf{nil}
    \EndIf
    \EndFunction
    \end{algorithmic}
\end{algorithm}

\paragraph{How to chose the upper limit?}  Our choice of keeping straight with
the limits of the standard is a mere choice of commodity: we are interested in
finding public keys  not respecting the standard.
Though, it is worth noting that what this limit \emph{states} is that at least
one of the most significant $100$ bits should be different between the two
primes:

\begin{bytefield}[
  endianness=big,
  bitwidth=1.35em,
  % bitformatting=\fakerange,
  ]{16}
  \\
  % \bitheader{}
  \\[1px]
  \begin{rightwordgroup}{\small{$2^{\frac{\log N}{2}-100}$}}
    \bitbox{1}{0} & \bitbox{1}{0} & \bitbox{1}{0} & \bitbox{1}{0} &
    \bitbox{1}{0} & \bitbox{1}{0} & \bitbox{1}{0} & \bitbox{1}{0} &
    \bitbox{3}{\tiny $\cdots$} &
    \bitbox{1}{0} & \bitbox{1}{0} & \bitbox{1}{0} & \bitbox{1}{0} &
    \bitbox{1}{0} & \bitbox{1}{1} & \bitbox{1}{0} & \bitbox{1}{0} &
    \bitbox{3}{\tiny $\cdots$}    & \bitbox{1}{0} & \bitbox{1}{0} &
  \end{rightwordgroup}
  \\[1ex]
  \wordbox[]{1}{} &&
  \\[1ex]
  \begin{rightwordgroup}{$p$}
    \bitbox{1}{0} & \bitbox{1}{1} & \bitbox{1}{0} & \bitbox{1}{0} &
    \bitbox{1}{0} & \bitbox{1}{0} & \bitbox{1}{1} & \bitbox{1}{1} &
    \bitbox{3}{\tiny $\cdots$} &
    \bitbox{1}{0} & \bitbox{1}{1} & \bitbox{1}{0} & \bitbox{1}{0} &
    \bitbox{1}{0} &
    \colorbitbox{lightgray}{1}{1} & \colorbitbox{lightgray}{1}{0} &
    \colorbitbox{lightgray}{1}{0} &
    \colorbitbox{lightgray}{4}{\tiny{$\cdots$ LSB $\cdots$}} &
    \colorbitbox{lightgray}{1}{0} &
  \end{rightwordgroup}
  \\[1ex]
  \begin{rightwordgroup}{$q$}
    \bitbox{1}{0} & \bitbox{1}{1} & \bitbox{1}{0} & \bitbox{1}{0} &
    \bitbox{1}{0} & \bitbox{1}{0} & \bitbox{1}{0} & \bitbox{1}{1} &
    \bitbox{3}{\tiny $\cdots$} &
    \bitbox{1}{0} & \bitbox{1}{1} & \bitbox{1}{0} & \bitbox{1}{0} &
    \bitbox{1}{0} &
    \colorbitbox{lightgray}{1}{0} & \colorbitbox{lightgray}{1}{0} &
    \colorbitbox{lightgray}{1}{0} &
    \colorbitbox{lightgray}{4}{\tiny{$\cdots$ LSB $\cdots$}} &
    \colorbitbox{lightgray}{1}{0} &
  \end{rightwordgroup}
\end{bytefield}
\vfill

For example, in the case of a RSA key $1024$, the binary difference between $p$
and $q$ has to be greater than $2^{412}$, which means that, excluding corner-cases
where the remainder is involved, there must be at least one difference in the
top 100 most significant bits for the key to be considered safe.


\section{Thoughts about a parallel solution}

At first glance we might be willing to split the entire interval
$\{ \ceil{\sqrt{N}}, \ldots, N-1 \}$ in equal parts, one per each
node. However, this would not be any more efficient than the trial division
algorithm, and nevertheless during each single iteration, the computational
complexity is dominated by the square root $\dsqrt$ function, which belongs to
the class \bigO{\log^2 N}, as we saw in section ~\ref{sec:preq:sqrt}.
Computing separatedly $x^2$ would add an overhead of the same order of magnitude
\bigO{\log^2 N}, and thus result in a complete waste of resources.

As a result of this, we advice the use of a strictly limited number of
processors - like two or three - performing in parallel fermat's factorization
method over different intervals.
%%% Local Variables:
%%% TeX-master: "question_authority.tex"
%%% End:

\chapter{Wiener's cryptanalysis method \label{chap:wiener}}

Wiener's attack was first published in 1989 as a result of cryptanalysis on the
use of short RSA secret keys ~\cite{wiener}. It exploited the fact that it is
possible to find the private key in \emph{polynomial time} using continued fractions
expansions whenever a good estimate of the fraction $\frac{e}{N}$ is known.
More specifically, given $d < \frac{1}{3} \sqrt[4]{N}$ one can efficiently
recover $d$ only knowing $\angular{N, e}$.

The scandalous implication behind Wiener's attack is that, even if there are
situations where having a small private exponent may be
particularly tempting with respect to performance (for example, a smart card
communication with a computer), they represent a threat to the security of the
cipher.
Fortunately, ~\cite{wiener} \S 6 presents a couple of precautions that make a
RSA key-pair immune to this attack, namely
(i) making $e > \sqrt{N}$ and
(ii) $gcd(p-1, q-1)$ large.

\section{A background on Continued Fractions \label{sec:wiener:cf}}

Let us call \emph{continued fraction} any expression of the form:
%% why \cfrac sucks this much. |-------------------------|
\begin{align*}
a_0 + \frac{1}{a_1
    + \frac{1}{a_2
    + \frac{1}{a_3
    + \frac{1}{a_4 + \ldots}}}}
\end{align*}
Consider now any \emph{finite continued fraction}, conveniently represented with
the sequence
$\angular{a_0, a_1, a_2, a_3,  \ \ldots, a_n}$.
Any number $x \in \mathbb{Q}$ can be represented as a finite continued fraction,
and for each $i < n$ there exists a fraction $\rfrac{h_i}{k_i}$ approximating
$x$.
By definition, each new approximation is recursively defined as:

\begin{align}
  \label{eq:wiener:cf}
  \begin{cases}
    a_{-1} = 0 \\
    a_i = h_i // k_i \\
  \end{cases}
  \quad
  \begin{cases}
    h_{-2} = 0 \\
    h_{-1} = 1 \\
    h_i = a_i h_{i-1} + h_{i-2}
  \end{cases}
  \quad
  \begin{cases}
    k_{-2} = 1 \\
    k_{-1} = 0  \\
    k_i = a_i k_{i-1} + k_{i-2}
  \end{cases}
\end{align}

Among the prolific properties of such objects, firstly Wiener ~\cite{wiener}
and later Boneh ~\cite{20years} discovered that, if a continued fraction $f'$ is
an underestimate of another one $f$, i.e.
\begin{align}
  f' = f(1-\delta)
\end{align}
then it is possible to recover $f$, having $f'$, if $\delta$ is ``small
enough'', where small enough means:
\begin{align}
  \label{eq:wiener:cf_approx}
  \delta = 1 - \frac{f'}{f} < \frac{1}{\rfrac{3}{2}{h_1}{k_1}}
\end{align}
\\
The \emph{continued fraction algorithm} allowing us to recover $f$ is the
following:
\begin{enumerate}[(i)]
  \setlength{\itemsep}{1pt}
  \setlength{\parskip}{0pt}
  \setlength{\parsep}{0pt}
  \item generate the next $a_i$ of the continued fraction expansion of $f'$;
  \item use ~\ref{eq:wiener:cf} to generate the next fraction $\rfrac{h_i}{k_i}$
    equal to $\angular{a_0, a_1, \ldots, a_{i-1}, a_i}$ %% non e` proprio cosi`
  \item check whether $\rfrac{h_i}{k_i}$ is equal to $f$
\end{enumerate}

\section{Constructing the attack}

As we saw in ~\ref{sec:preq:rsa}, by construction the two exponents are such that
$ed \equiv 1 \pmod{\varphi(N)}$. This implies that there exists a
$k \in \naturalN \mid ed = k\varphi(N) + 1$. This can be formalized to be
the same problem we formalized in ~\ref{sec:wiener:cf}:
\begin{align*}
  ed = k\varphi(N) + 1 \\
  \abs{\frac{ed - k\eulerphi{N}}{d\eulerphi{N}}} = \frac{1}{d\eulerphi{N}} \\
  \abs{\frac{e}{\eulerphi{N}} - \frac{k}{d}} = \frac{1}{d\eulerphi{N}} \\
\end{align*}

Now we proceed by substituting $\eulerphi{N}$ with $N$, since for large $N$, one
approximates the other. We consider also the difference of the two, limited by
$\abs{\cancel{N} + p + q - 1 - \cancel{N}} < 3\sqrt{N}$.
For the last step, remember that $k < d < \rfrac{1}{3}\sqrt[4]{N}$:

\begin{align*}
  \abs{\frac{e}{N} - \frac{k}{d}} &= \abs{\frac{ed - kN}{Nd}} \\
  &= \abs{\frac{\cancel{ed} -kN - \cancel{k\eulerphi{N}} + k\eulerphi{N}}{Nd}} \\
  &= \abs{\frac{1-k(N-\eulerphi{N})}{Nd}} \\
  &\leq \abs{\frac{3k\sqrt{N}}{Nd}}
  = \frac{3k}{d\sqrt{N}}
  < \frac{3(\rfrac{1}{3}\ \sqrt[4]{N})}{d\sqrt{N}}
  = \frac{1}{d\sqrt[4]{N}} < \frac{1}{2d^2}
\end{align*}

This demonstrates the conditions of ~\ref{eq:wiener:cf_approx} holds, and allows
us to proceed with the continued fraction algorithm to converge to a solution
~\cite{20years}.

\paragraph{}
We start by generating the $\log N$ continued fraction expansions of
$\frac{e}{N}$, and for each convergent $\frac{k}{d}$,
%% XXX. verify this
which by contruction is already at the lowest terms, we verify if it produces a
factorization of $N$.
First we check that $\eulerphi{N} = \frac{ed-1}{k}$ is
an integer. Then we solve ~\ref{eq:wiener:pq} in $x$ in order to find $p, q$:
\begin{align}
  \label{eq:wiener:pq}
  x^2 - (N - \eulerphi{N} + 1)x + N = 0
\end{align}
The above equation is constructed so that the $x$ coefficient is the sum of the
two primes, while the constant term $N$ is the product of the two. Therefore, if
$\eulerphi{N}$ has been correctly guessed, the two roots will be $p$ and $q$.

\section{Again on the engine™}

The algorithm is pretty straightforward by itself: we just need to apply the
definitions provided in ~\ref{eq:wiener:cf} and test each convergent until
$\log N$ iterations have been reached.
%% XXX. questo viene da 20 years, ma non e` spiegato perche`.
A Continued fraction structure may look like this:

\begin{minted}{c}
  typedef struct cf {
    bigfraction_t fs[3];  /* holding h_i/k_i, h_i-1/k_i-1, h_i-2/k_i-2 */
    short i;              /* cycling in range(0, 3) */
    bigfraction_t x;      /* pointer to the i-th fraction in fs */
    BIGNUM* a;            /* current a_i */
    BN_CTX* ctx;
  } cf_t;
\end{minted}
where \texttt{bigfraction\_t} is just a pair of \texttt{BIGNUM} \!s
$\angular{h_i, k_i}$. Whenever we need to produce a new convergent, we increment
$i \pmod{3}$ and apply the definitions given. The fresh convergent must be
tested with very simple algebraic operations. It is worth noting here that
\ref{eq:wiener:pq} can be solved using the reduced discriminant formula, as
$p, q$ are odd primes:
\begin{align*}
\Delta = \left( \frac{N-\eulerphi{N} + 1}{2} \right)^2 - N \\
x_{\angular{p , q}} = - \frac{N - \eulerphi{N} + 1}{2} \pm \sqrt{\Delta}
\end{align*}
Assuming the existence of the procedures \texttt{cf\_init}, initializing a
continued fraction structure, and \texttt{cf\_next} producing the next
convergent, we provide an algorithm for attacking the RSA cipher via Wiener:

\begin{algorithm}[H]
  \caption{Wiener's Attack}
  \label{alg:wiener}
  \begin{algorithmic}[1]
    \Function{wiener}{\PKArg}
    \State $f \gets  \texttt{cf\_init}(e, N)$
    \For{$\ceil{\log N} \strong{ times }$}
      \State $k, d \gets \texttt{cf\_next}(f)$
      \If{$k \nmid ed-1$} \strong{continue} \EndIf
      \State $\eulerphi{N} \gets (ed - 1)\ //\ k$
      \If{$\eulerphi{N}$ is odd} \strong{continue} \EndIf
%% XXX. it could be that calling 'b' b/2 and 'delta' sqrt(delta/4) is
%% misleading.
      \State $b \gets (N - \eulerphi{N} + 1) \gg 1$
      \State $\Delta, r \gets \dsqrt{b^2 - N}$
      \If{$r \neq 0$} \strong{continue} \EndIf
      \State $p \gets b + \Delta$
      \State $q \gets b - \Delta$
      \State \strong{break}
    \EndFor
    \State \Return p, q
    \EndFunction
  \end{algorithmic}
\end{algorithm}

\section{Building a distributed version}
%%% Local Variables:
%%% mode: latex
%%% TeX-master: "question_authority"
%%% End:

\include{pollard-1}
\chapter{Pollard's $p+1$ factorization method}

pollard!
%%% Local Variables:
%%% mode: latex
%%% TeX-master: "question_authority"
%%% End:

\chapter{Pollard's $\rho$ factorization method \label{chap:pollardrho}}

A \emph{Monte Carlo} factorization method, published by J. M. Pollard in
~\cite{pollardMC}, consists into identifying a periodically recurrent  sequence
of integers within a random walk $\pmod{N}$ that could leak one of the two
factors.

Consider a function $f$ from $\mathcal{S}$ to $\mathcal{S}$, where
$\mathcal{S} = \{0, 1, \ldots, q-1\}$ and $q \in \naturalPrime$.
Let $s$ be a random element in $\mathcal{S}$, and consider the sequence
\begin{align*}
  s,\ f(s),\ f(f(s)),\ \ldots
\end{align*}
Since $f$ acts over a finite set, it is clear that this sequence must
eventually repeat, and become cyclic.
We might diagram it with the letter $\rho$, where the tail represent the
aperiodic part, or \emph{epacts}, and the oval the cyclic part, or
\emph{period}.
\begin{center}
  \begin{tikzpicture}[scale=0.7, thick]
    \tikzstyle{every node}=[draw,circle,fill=white,minimum size=4pt,
                            inner sep=0pt]
    \node (1) at (1.4, 0.2) [label=left:$x_1$] {};
    \node (2) at (2.5, 3)   [label=left:$x_{i-2}$] {};
    \node (3) at (3.25, 5)  [label=left:$x_{i-1}$] {};
    \node (4) at (4, 7)     [label=left:$ x_i \equiv x_j $] {};
    \node (5) at (6, 9)     [label=above:$x_{i+1}$] {};
    \node (6) at (8, 7)     [label=right:$x_{i+2}$] {};
    \node (7) at (6, 5)     [label=below:$x_{j-1}$] {};

    \path (1) edge [dashed] (2);
    \path (2) edge (3);
    \path (3) edge (4);
    \path (4) edge [bend left] (5);
    \path (5) edge [bend left] (6);
    \path (6) edge [bend left, dashed] (7);
    \path (7) edge [bend left] (4);

    %%\draw [decorate,decoration={brace, raise=1.5cm}] (1) -- (3)
    %%node[draw=no] at (-1.5, 4) {tail};
    \draw [decorate,decoration={brace, raise=3cm}] (5) -- (7)
    node[draw=none] at (13, 7) {\footnotesize {period $\pmod{q}$}};

\end{tikzpicture}
\end{center}

Now, consider $N=pq$.
Let $F(x)$ be any function generating pseudorandom integers
$\angular{x_1, x_2, \ldots}$, and let $f(x) = F(x) \pmod{q}$.
As we said above, without any luck, there will be a pair $\angular{x_i, x_j}$
generated by $F$ such that $x_i \equiv x_j \pmod{q}$, but $x_i \neq x_j$.

Therefore, in order to factorize $N$, we proceed as follows: starting from a
random $s$, we iteratively apply $F$ reduced modulo $N$. Whenever we find a
period, if $\gcd(x_i - x_j, N) > 1$ then we found a non-trivial
factor of $N$.

\paragraph{Choosing the function} Ideally, $F$ should be easily computable, but
at the same time random enough to reduce as much as possible the epacts
~\cite{Crandall} \S 5.2.1. Any quadratic function $F(x) = x^2 + b$ should be
enough\footnote{
  Note that this has been only empirically verified, and so far not been proved
  (~\cite{riesel}, p. 177)},
provided that $b \in \naturalN \setminus \{0, 2\}$.
For example, ~\cite{pollardMC} uses $x^2 -1$, meanwhile we are going to choose
$F(x) = x^2 + 1$.

\paragraph{Finding the period} The trivial way to discover a period would be to
test $x_i$ with all $x_j, \quad j < i$. Though, in \cite{AOCPv2} \S 3.1,
Knuth gives a simple and elegant algorithm, attributed to Floyd, for finding a
multiple of the period.
This algorithm is the same one finally adopted by Pollard in
~\cite{pollardMC}.

\begin{theorem*}[Floyd]
Given an \emph{ultimately periodic} sequence, in the sense that there exists
numbers $\lambda$ and $\mu$ for which the values:
\begin{align*}
  X_0, X_1, \ldots, X_{\mu}, \ldots, X_{\mu + \lambda - 1}
\end{align*}
are distinct, but $X_{n+\lambda} = X_n$ when $n \geq \mu$,
then there exists an
$\mu < n < \mu + \lambda$ such that $X_n = X_{2n}$.
\end{theorem*}

\begin{proof}
  First, if $X_n = X_{2n}$, then the sequence is obviously periodic from
  $X_{2n}$ onward, possibly even earlier.
  Conversely, it is true that $X_n = X_m \quad (n \geq \mu)$ for
  $m = n + k\lambda, \quad k \in \naturalN$. Hence, there will eventually
  be an $n$ such that $X_n = X_{2n}$ if and only if $n - \mu$ is a multiple of
  $\lambda$.
  The first such value happens for $n = (\lambda + 1)\floor{\rfrac{\mu}{\lambda}}$.
\end{proof}

The immediate consequence of this is that we can find a collision simply by
checking $\gcd(x_{2i} - x_i, N) > 1$ for incremental $i$.

\paragraph{Brent's Improvement} In 1979, Brent discovered an entire family of
cycle-finding algorithms whose optimal version resulted to be 36\% faster than
Floyd's one \cite{pollard-brent}.
Instead of looking for the period of the sequence using $x_{2i} - x_i$, Brent
considers
$\abs{x_j - x_{2^k}}$ for $ 3 \cdot 2^{k-1} < j \leq 2^{k+1}$, resulting in
fewer operations required by the algorithm. Pragmatically, this boils down to
compare:

\medskip
\begin{tabular}{l@{\hskip 40pt} l@{\hskip 50pt} l}
  $k = 0$ & $j \in \{1+1 \ldots 2\}$ & $\abs{x_1 - x_2}$ \\
  $k = 1$ & $j \in \{3+1 \ldots 4\}$ & $\abs{x_2 - x_4}$ \\
  $k = 2$ & $j \in \{6+1 \ldots 8\}$ & $\abs{x_4 - x_7}$, $\abs{x_4 - x_8}$ \\
  $k = 3$ & $j \in \{12+1, \ldots 16\}$ &
            $\abs{x_8 - x_{13}}$, $\abs{x_8 -x_{14}}$, $\ldots$, $|x_8 - x_{16}|$\\
  $k = 4$ & $j \in \{24+1, \ldots 32 \}$ &
            $\abs{x_{16} - x_{25}}$, $\ldots$, $\abs{x_{16} - x_{32}}$ \\[2pt]
  $\quad \vdots$ & $\quad \vdots$ & $\quad \quad \vdots$ \\
  %\multicolumn{1}{c}{$\vdots$} &
  %\multicolumn{1}{c}{$\vdots$} &
  %\multicolumn{1}{c}{$\vdots$} \\
\end{tabular}


A Pollard's $\rho$ variant that implements Brent's cycle-finding algorithm
instead of Floyd's one runs around 25\% faster on average
~\cite{pollard-brent}.

\section{Complexity}
\cite{riesel} presents a nice proof of the \emph{average} complexity of
this algorithm, based on the birthday paradox.
\newtheorem*{birthday}{The Birthday Paradox}
\begin{birthday}
  How many persons needs to be selected at random in order that the probability
  of at least two of them having the same birthday exceeds $\rfrac{1}{2}$?
\end{birthday}

\begin{proof}[Solution]
  The probability that $\epsilon$ different persons have different birthdays is:
  \begin{align*}
    \Big(1 - \frac{1}{365}\Big)
    \Big(1 - \frac{2}{365}\Big)
    \Big(1 - \frac{3}{365}\Big)
    \cdots
    \Big(1 - \frac{\epsilon -1}{365}\Big)
    =
    \frac{365!}{365^\epsilon (365-\epsilon)!}
  \end{align*}
  This expression becomes $< \rfrac{1}{2}$ for $\epsilon \geq 23$.
\end{proof}

We can obviously substitute the $365$ with any set of cardinality $\zeta$
to express the probability that a random function from $\integerZ_{\epsilon}$
to $\integerZ_{\zeta}$ is injective. However, back to our particular case,
we want to answer the question:

\emph{
  How many random numbers do we have to run through before finding at least
  two integers equivalent $\mod{q}$?
}


Using the same reasoning presented above over the previously defined function
$f(x): \mathcal{S} \to \mathcal{S}$, we will discover that after
$\approx 1.18 \sqrt{q}$ steps the probability to have fallen inside the period
is $\rfrac{1}{2}$. %% is it clear that q is either one of the two primes, and
                   %% that here we want to examinate only a portion of the
                   %% domain?
Since any of the two primes factoring $N$ is bounded above by $\sqrt{N}$, we
will find a periodic sequence, and thus a factor, in time \bigO{\sqrt[4]{N}}.


\section{An Implementation Perspective}

The initial algorithm described by Pollard \cite{pollardMC} and consultable
immediately below, looks for the pair $\angular{x_i, x_{2i}}$ such that
$\gcd(x_{2i} - x_i, N) > 1$.  This is achieved by keeping two variables $x, y$
and respectively updating them via $x \gets f(x)$ and $y \gets f(f(y))$.

\begin{algorithm}
  \caption{Pollard's $\rho$ factorization}
  \begin{algorithmic}[1]
    \Function{rho}{\PKArg}
    \State $x \getsRandom \naturalN$
    \State $y \gets x$
    \State $g \gets 1$
    \While{$g = 1$}
      \State $x \gets x^2 + 1 \pmod{N}$
      \State $y \gets y^4 + 2y^2 + 2 \pmod{N}$
      \State $g \gets gcd(|x - y|, N)$
    \EndWhile
    \If{$g = N$} \Return \strong{nil}
    \Else \ \ \Return $g$
    \EndIf
    \EndFunction
  \end{algorithmic}
\end{algorithm}

\begin{remark}
  It is intresting to see how in its basic version, Pollard's $\rho$
  method just needs 3 variables  to preserve the
  state. This places it among the most parsimonious factorization algorithms in
  terms of memory footprint.
\end{remark}

An immediate improvement of this algorithm would be to occasionally compute Euclid's
algorithm over the accumulated product to save some computation cycles, just as
we saw in section~\ref{sec:pollard-1:implementing}. The next code fragment
adopts this trick together with Brent's cycle-finding variant:

\paragraph{Parallelism}
Unfortunately, a parallel implementation of the $\rho$ algorithm would not
provide a linear speedup.

The computation of the $x_i$ sequence is intrinsecally serial; the only
plausible approach to parallelism would be to try several different pseudorandom
sequences, in which case $m$ different machines processing $m$ different
sequences in parallel would be no more than \bigO{\sqrt{m}}
efficient (\cite{brent:parallel} \S 3).

\begin{algorithm}
  \caption{Pollard-Brent's factorization \label{alg:pollardrho}}
  \begin{algorithmic}[1]
    \Function{rho}{\PKArg}
    \State $s \gets 100$
    \Comment steps to check for $\gcd$
    \State $i \gets 1; \quad j' \gets j \gets 1$
    \Comment step and period counters
    \State $x' \gets x \getsRandom \naturalN$
    \Comment the $x_i$ discussed above
    \State $y' \gets y \gets x^2 + 1$
    \Comment the $x_j$ discussed above
    \State $k \gets 0; \quad q \gets \abs{x-y}$
    \While{$g = 1$}
      \State $x \gets y$ \Comment $x_i = 2^k$
      \State $j \gets 3 \cdot 2^{k} + 1$
      \State $k \gets k+1$
      \While{$j \leq 2^k$}
        \State $j \gets j+1; \quad i \gets i+1$
        \Comment increment step and period counters
        \State $y \gets y^2 + 1 \pmod{N}$
        \State $q \gets q \cdot \abs{x - y}$
        \If{$ i \mid s$} \Comment time to compute $\gcd$?
          \State $g \gets \gcd(q, N)$
          \If{$g = N$}
            \Comment too far: fall back to latest epoch
            \State $s \gets 1; \quad  g \gets 1$
            \State $j \gets j'; \quad x \gets x'; \quad y \gets y'$
          \ElsIf{$g = 1$}
            \Comment save current state
            \State $x' \gets x; \quad y' \gets y$
            \State $j' \gets j$
          \Else
            \State $p \gets g$
            \State $q \gets N//p$
            \State \Return $p, q$
          \EndIf
        \EndIf
      \EndWhile
    \EndWhile
    \EndFunction
  \end{algorithmic}
\end{algorithm}

%%% Local Variables:
%%% mode: latex
%%% TeX-master: "question_authority"
%%% End:

\chapter{Dixon {\texttt{\small{[insert random buzzword here]}}} method\label{chap:dixon}}

~\cite{dixon} describes a class of ``probabilistic algorithms'' for finding a
factor of any composite number, at a sub-exponential cost. They basically
consists into taking random integers $r$ in $\{1, \ldots, N\}$ and look for those
where $r^2 \mod{N}$ is \emph{smooth}. If enough are found, then those integers
can somehow be assembled, and so a fatorization of N attemped.

%% that's not really academic to be stated officially, but I would have never
%% understood this section without Firas (thanks).
%% <http://blog.fkraiem.org/2013/12/08/factoring-integers-dixons-algorithm/>
%% I kept the voila` phrase, that was so lovely.
\section{A little bit of History}
During the latest century there has been a huge effort to approach the problem
formulated by Fermat ~\ref{eq:fermat_problem} from different perspecives. This
led to an entire family of algorithms known as \emph{Quadratic Sieve} [QS]. The
core idea is still to find a pair of perfect squares whose difference can
factorize $N$, but maybe Fermat's hypotesis can be made weaker.

\paragraph{Kraitchick} was the first one popularizing the idea the instead of
looking for integers $\angular{x, y}$ such that $x^2 -y^2 = N$ it is sufficient
to look for \emph{multiples} of $N$:
\begin{align}
  x^2 - y^2 \equiv 0 \pmod{N}
\end{align}
and, once found, claim that $\gcd(N, x \pm y)$ are non-trial divisors of $N$
just as we did in \ref{sec:fermat:implementation}.
Kraitchick did not stop here: instead of trying $x^2 \equiv y^2 \pmod{N}$ he
kept the value of previous attempt, and tries to find \emph{a product} of such
values which is also a square. So we have a sequence
\begin{align}
  \label{eq:dixon:x_sequence}
  \angular{x_0, \ldots, x_k} \mid \forall i \leq k \quad x_i^2 - N
  \; \text{ is a perfect square}
\end{align}
and hence
\begin{align*}
  \prod_i (x_i^2 - N) = y^2
\end{align*}
that $\mod{N}$ is equivalent to:
\begin{align}
  \label{eq:dixon:fermat_revisited}
  y^2 \equiv \prod_i x_i^2 - N \equiv \big( \prod_i x_i \big) ^2 \pmod{N}
\end{align}
and voil\`a our congruence of squares. For what concerns the generation of $x_i$
with the property \ref{eq:dixon:x_sequence}, they can simply taken at random and
tested using trial division.

\paragraph{Brillhart and Morrison} later proposed (\cite{morrison-brillhart}
p.187) a better approach than trial division to find such $x$. Their idea aims
to ease the enormous effort required by the trial division. In order to achieve
this. they introduce a \emph{factor base} $\factorBase$ and generate random $x$
such that $x^2 - N$ is $\factorBase$-smooth. Recalling what we anticipated in
~\ref{sec:preq:numbertheory}, $\factorBase$ is a precomputed set of primes
$p_i \in \naturalPrime$.
This way the complexity of generating a new $x$ is dominated by
\bigO{|\factorBase|}. Now that the right side of \ref{eq:dixon:fermat_revisited}
has been satisfied, we have to select a subset of those $x$ so that their
product can be seen as a square. Consider an \emph{exponent vector}
$v_i = (\alpha_0, \alpha_1, \ldots, \alpha_r)$ associated with each $x_i$, where
\begin{align*}
  a_j = \begin{cases}
    1 \quad \text{if $p_j$ divides $x_i$ to an odd power} \\
    0 \quad \text{otherwise}
    \end{cases}
\end{align*}
for each $0 \leq j \leq r $. There is no need to restrict ourselves for positive
values of $x^2 -N$, so we are going to use $\alpha_0$ to indicate the sign. This
benefit has a neglegible cost: we have to add the non-prime $-1$ to our factor
base.

Let now $\mathcal{M}$ be the rectangular matrix having per each $i$-th row the
$v_i$ associated to $x_i$: this way each element $m_{ij}$ will be $v_i$'s
$\alpha_j$. We are interested in finding set(s) of $x$ that satisfies
\ref{eq:dixon:fermat_revisited}, possibly all of them.
Define $K$ as the subsequence of $x_i$ whose product always have even powers.
This is equivalent to look for the set of vectors $\{ w \mid wM = 0 \}$ by
definition of matrix multiplication in $\mathbb{F}_2$.


\paragraph{Dixon} Morrison and Brillhart's ideas of \cite{morrison-brillhart}
were actually used for a slightly different factorization method, employing
continued fractions instead of the square difference polynomial. Dixon refined
those by porting to the quare problem, achieving a probabilistic factorization
method working at a computational cost asymptotically  best than all other ones
previously described: \bigO{\beta(\log N \log \log N)^{\rfrac{1}{2}}} for some
constant $\beta > 0$ \cite{dixon}.

\section{Computing the Kernel}

%%% Local Variables:
%%% mode: latex
%%% TeX-master: "question_authority"
%%% End:


\chapter{Epilogue}

This project led to this this this that
%%% Local Variables:
%%% mode: latex
%%% TeX-master: "question_authority"
%%% End:


\backmatter
%%\bibliographystyle{ieeetr}
\bibliography{library}
\clearpage
\addcontentsline{toc}{chapter}{Bibliography}

\end{document}
